% Bagian: incompleteness
% Bab: theories-computability
% Subbagian: extensions-of-q-not-decidable

\documentclass[../../../include/open-logic-section]{subfiles}

\begin{document}

\olfileid[id]{inc}{tcp}{cqn}

\olsection{Perluasan $\Th{Q}$ yang Konsisten Tidak Dapat Diputuskan}

\begin{explain}
Ingat bahwa suatu teori \emph{konsisten} jika teori itu tidak membuktikan
$!A$ sekaligus $\lnot !A$ untuk sebarang formula~$!A$. Karena apa pun
merupakan konsekuensi logis dari suatu kontradiksi, teori yang tak konsisten bersifat trivial:
setiap kalimat dapat dibuktikan. Jelas bahwa jika suatu teori
$\omega$-konsisten, maka teori itu konsisten. Namun, konsistensi merupakan
syarat yang lebih lemah (yakni terdapat teori yang konsisten tetapi tidak
$\omega$-konsisten). Kita dapat memperlemah asumsi dalam
\olref[oqn]{thm:oconsis-q} menjadi konsistensi biasa untuk memperoleh teorema
yang lebih kuat.
\end{explain}

\begin{lem}
Tidak ada ``relasi dapat dihitung universal.'' Artinya, tidak ada relasi
dapat dihitung dua-tempat~$R(x,y)$ dengan sifat berikut: setiap kali $S(y)$
merupakan relasi dapat dihitung satu-tempat, terdapat suatu~$k$ sedemikian
sehingga, untuk setiap~$y$, $S(y)$ benar jika dan hanya jika $R(k,y)$ benar.
\end{lem}

\begin{proof}
Andaikan $R(x,y)$ merupakan relasi dapat dihitung universal. Definisikan
$S(y)$ sebagai relasi $\lnot R(y,y)$. Karena $S(y)$ dapat dihitung, untuk suatu~$k$,
$S(y)$ ekuivalen dengan~$R(k,y)$. Namun, dengan demikian $S(k)$ ekuivalen
dengan $R(k,k)$ sekaligus $\lnot R(k,k)$, dan ini merupakan kontradiksi.
\end{proof}


\begin{thm}
Misalkan $\Th{T}$ sebarang teori konsisten yang memuat~$\Th{Q}$. Maka
$\Th{T}$ tidak dapat diputuskan.
\end{thm}


\begin{proof}
Andaikan $\Th{T}$ merupakan perluasan $\Th{Q}$ yang konsisten dan dapat
diputuskan. Kita akan memperoleh kontradiksi dengan menggunakan~$\Th{T}$
untuk mendefinisikan relasi dapat dihitung universal.

Misalkan $R(x,y)$ berlaku jika dan hanya jika
\begin{quote}
$x$ mengodekan suatu formula~$!D(u)$, dan $\Th{T}$ membuktikan
$!D(\num y)$.
\end{quote}
Karena kita mengasumsikan bahwa $\Th{T}$ dapat diputuskan, $R$ dapat dihitung.
Mari kita tunjukkan bahwa $R$ bersifat universal. Jika $S(y)$ sebarang relasi yang
dapat dihitung, maka relasi itu dapat direpresentasikan dalam~$\Th{Q}$ (dan
karena itu dalam~$\Th{T}$) oleh suatu formula~$!D_S(u)$. Maka untuk setiap~$n$,
kita mempunyai
\begin{eqnarray*}
S(\num n) & \lif & T \vdash !D_S(\num n) \\
& \lif & R(\Gn{!D_S(u)}, n)
\end{eqnarray*}
dan
\begin{eqnarray*}
\lnot S(\num n) & \lif & T \vdash \lnot !D_S(\num n) \\
& \lif & T \not\vdash !D_S(\num n) \quad \text{(karena $\Th{T}$
  konsisten)} \\
& \lif & \lnot R(\Gn{!D_S(u)}, n).
\end{eqnarray*}
Artinya, untuk setiap~$y$, $S(y)$ benar jika dan hanya jika
$R(\Gn{!D_S(u)}, y)$ benar. Jadi $R$ bersifat universal, dan kita memperoleh
kontradiksi yang dicari.
\end{proof}

Misalkan ``aritmetika sejati'' sebagai teori $\Setabs{!A}{ \Sat{\Nat}{!A}}$,
yakni himpunan kalimat dalam bahasa aritmetika yang benar dalam interpretasi
standar.

\begin{cor}
Aritmetika sejati tidak dapat diputuskan.
\end{cor}

%Dalam Bagian~\olref{undefinability:truth:section} kita akan menyatakan hasil
%yang lebih kuat, yang disebabkan oleh Tarski.

\end{document}
